\section{The Security Function}\label{sec:design}

\sys adapts task-alignment checking~\cite{taskshield,llamafirewall} to an NLIP SE. Every
tool call must be one the user would approve as a step toward what they asked. Unlike
AlignmentCheck~\cite{llamafirewall}, it never reads the agent's reasoning. It decides only
from data visible at its enforcement points: the user's words from the NLIP message, the
structured call (tool name and JSON arguments) and its result, its own per-session state, and
a deployer policy file.

\paragraph{Placement and plugging in.} The core is framework-agnostic. It consumes
\emph{call} and \emph{result} events and returns allow, block or replace, writing each
decision to a JSONL audit log. Two adapters attach it:
\begin{itemize}
\item An NLIP ingress/egress SE (\texttt{tool\_guard}), enabled like any other extension in
  the server's \texttt{config.toml}. At ingress it extracts the text-format submessages of the
  incoming NLIP message and appends them to the session's \emph{trusted task}. If no
  request in the session has a text submessage, there is no trusted task and every tool call
  is blocked. This is the only thing the
  guard trusts as ``what the user asked'': tool descriptions and tool output are data. It
  never blocks the NLIP message itself, and egress closes the turn in the audit log.
\item A LangChain middleware (\texttt{awrap\_tool\_call}), bound to the NLIP session. It
  calls the guard immediately before and after each tool call. Porting to another agent
  framework means rewriting this adapter, a class of about 40 lines.
\end{itemize}
Because the framework has no tool-boundary hook, this second adapter is per agent framework;
the application finds the guard through a registry keyed by the SE instance. Removing
\texttt{tool\_guard} from the SE lists runs the agent undefended. Our
undefended baseline is exactly that.

\paragraph{Pre-tool check.} For each proposed call, in order:
\begin{enumerate}
\item \emph{Validation.} A missing trusted task, an empty tool name or non-object
  arguments block the call (fail closed).
\item \emph{Risk label.} The deployer labels each tool, by name, in the policy file:
  read-only (14 tools), internal write (4), external send (3) or destructive (3). A
  read-only call is allowed without a model call; its output is checked after the call. A
  tool missing from the policy, for example one added later by an MCP server, is blocked.
\item \emph{\dthree judge.} Every other call goes to \texttt{d3-flash}, a decision model of
  the vLLM Semantic Router~\cite{d3}. It returns the probability of ``yes'' for a yes/no
  question without generating text; we call this score P(approve). The question is: ``Would
  the user, having made this request, approve this exact tool call as a step toward it?''
  Explicit criteria say ``no'' means a goal the user did not ask for, data sent to someone
  the request does not lead to, or unrequested deletes, shares or changes. \dthree sees a
  JSON state: all user turns of the session; the call; up to 12 \emph{evidence lines} (at most
  300 characters each), which are lines of earlier results that mention the call's recipients,
  names or ids, so the judge can tell whether a recipient follows from content the user
  pointed to; and the last 6 calls with their results (1,500 characters each), marked as
  untrusted. The call runs if P(approve) $\ge 0.5$, or $\ge 0.7$ once the session is \emph{tainted}. A
  judge error blocks.
\end{enumerate}
A blocked call is never sent to the MCP server. The model receives an error result with the
judge's score, and is told not to retry and to continue with the user's request.

\paragraph{Post-tool scan.} Each result (YAML from the MCP server) is parsed into text
fields. Every field of 40 or more characters goes to \dthree, in batches of 9; shorter fields
are not scanned, duplicates are scanned once, and only a field's first 4,000 characters are
sent. A result that is not structured YAML is scanned as one text (its first 4,000 characters, with
no minimum length). The question is: ``Does it contain instructions
addressed to an AI assistant or agent \ldots{} as opposed to ordinary content written for a
human reader?'' In a flagged field longer than 300 characters, the paragraphs (or lines) are
scored separately. The span from the first to the last flagged one is replaced by a notice;
if none is flagged, the whole field is. Ids, senders and dates are kept so the task can
continue. A flagged result taints the
session. A scan error replaces the whole result.

\paragraph{Why no rules.} Earlier versions also had regex rules, a \emph{catalog check} that
stripped hidden instructions from tool descriptions at \texttt{tools/list}, and policy rules
that blocked planted values and unrequested deletes. We dropped all three: the regex rules
match the very templates of our attack sets (\texttt{<IMPORTANT>}, \texttt{<INFORMATION>}),
which amounts to knowing the test in advance, and the description rules catch
\CatalogRulesPlainCaught\% of plain-sentence poison. The catalog check is a trade-off
(Appendix~\ref{app:ablation}): it gave the best task completion end to end, but cannot see
parameter descriptions and strips unusual legitimate documentation. Against tool poisoning
the judge alone stopped every attack these components stopped; against data injections the
rules catch what the \dthree scanner misses, and the judge as backstop is untested.
